\section{Considerações finais}
\label{sec:consideracoes}

A pergunta central foi respondida pela integração de duas camadas parcialmente sobrepostas. O mapeamento bibliométrico identificou transição de eficiência, desempenho e edificações verdes para BIM, \textit{digital twins}, carbono, ciclo de vida e manutenção preditiva. A síntese temática mostrou que dimensões técnica-operacional, institucional e ambiental e critérios de desempenho, informação, custo, vida útil, energia e risco formam o núcleo documental. Dos 121 registros temáticos, 109 também pertencem ao estrato bibliométrico de 372. Essa integração fundamentou a matriz de critérios e indicadores e o fluxo de validação para futura parametrização multicritério.

Foram identificadas seis dimensões, 15 critérios e abordagens decisórias heterogêneas. Estruturas genéricas, BIM e apoio à decisão foram mais frequentes que métodos multicritério nomeados. A busca de sensibilidade elevou o aprendizado de máquina de nove para 26 registros. Entre os 17 incorporados, dez concentram-se em previsão ou previsão com otimização, dois em diagnóstico e classificação e cinco em síntese ou integração tecnológica. Em conjunto, eles cobrem partes da cadeia entre modelo treinado, validação externa, ponderação, vetos e decisão pública auditável.

Edificações genéricas e portfólios predominaram, enquanto universidades, \textit{campus} e edifícios públicos foram menos frequentes. Doze registros apresentaram lacunas específicas para instituições públicas de ensino superior, resultado detalhado na Seção~\ref{sec:aplicabilidade}. Os ODS apareceram em um registro, e os componentes ambientais, sociais e institucionais foram mais frequentes que o rótulo ESG \parencite{masmoudi_ahptopsis_2026,ramantswana_esgfm_2026}.

A contribuição central é uma especificação operacional candidata e rastreável. A revisão sustenta os critérios; indicadores, unidades, fontes, periodicidades e direções são proposições autorais para validação. A definição de pesos, agregação, limiares e vetos integra a etapa institucional seguinte. Quatro frentes organizam essa continuidade:

\begin{enumerate}
\item \textbf{Validação de conteúdo:} painel de manutenção, engenharia, sustentabilidade, planejamento, orçamento e usuários revisa dimensões, indicadores e vetos antes da elicitação dos pesos.
\item \textbf{Integração de dados:} ordens de serviço, inspeções, custos, consumo, riscos e reclamações recebem dicionário, periodicidade e responsáveis.
\item \textbf{Construção de protótipo:} método de elicitação, normalização, pesos, agregação, limiares, casos retrospectivos e análise de sensibilidade são parametrizados e testados.
\item \textbf{Validação multicampi:} estabilidade, utilidade decisória e transferência são comparadas entre unidades e outros portfólios públicos.
\end{enumerate}

Diante da escassez de recursos, esta especificação oferece um caminho para justificar cada investimento por critérios documentados e auditáveis, convertendo a priorização da manutenção de resposta improvisada em prática de governança e valor público.